\section{Introduction}
\label{sec:introduction}

\oracle{} is the data acquisition for collinear laser spectroscopy (CLS) at IGISOL. It steps
the platform voltage through a scan, records every ion the time tagger sees against the
voltage it was taken at, and calibrates the voltage axis before every run. It keeps the
collinear DAQ's layout, file format and functions, brings in the bench calibration tool
cls-cal~\cite{clscal}, and is built so that the faults found in the old DAQ cannot happen
rather than being caught afterwards (Appendix~\ref{sec:design}). Runs are read by
cls\_tools~\cite{clstools} and fitted in \denis{}~\cite{denis}, both unchanged.

The window has five tabs above the quick-tune controls and the experiment information:

\begin{description}
    \item[\app{Scan controls}.] Scan and buncher parameters, run control, the scan queue,
        the ELOG panel and the live plots (Section~\ref{sec:scan}).
    \item[\app{Scan Calibration}.] A stand-alone calibration and the gate's verdict on it;
        every scan's own calibration is shown here too (Section~\ref{sec:calibration}).
    \item[\app{Hardware Calibration}.] The bench tool's measurements of the rack itself:
        settle time, linearity, the cooler, the bench table and the health check
        (Section~\ref{sec:hardware}).
    \item[\app{Scan Performance}.] How the rack's behaviour during a run shaped its
        spectrum, check by check, with the run's systematic inputs
        (Section~\ref{sec:performance}).
    \item[\app{Time graphs}.] Counts against time, to watch the beam.
\end{description}

\paragraph{Principles.}
Every calibration point is written, confirmed and only then read, so the voltage axis
cannot slip. Nothing is cut: what deserves a look is flagged, and the operator decides in
\denis{}. Every input survives a crash, and a run can be reopened exactly as it was taken.

\paragraph{Status.}
\oracle{} runs end to end on its simulated rack, which imitates the rack the bench tool
measured in October 2026. It has not yet driven the beamline: its drivers follow code that
has run on the rack but are tested against stand-ins for the vendor libraries.

\paragraph{Scope.}
This is a user manual. The code (\code{cls\_acquisition/}, \code{gui/}) documents itself
for developers; the README describes the test suite.
